% Plantilla de presentación - PFM01 / IIIEPE
% Alumno: Mtro. Martin Jonathan de la Cruz Muñoz
% Matrícula: 00874
%
% ==============================================================================
% PLANTILLA BASE DE ACTIVIDADES - PFM01
% Fundamentos para la enseñanza y el aprendizaje
% Maestría en Aprendizaje y Enseñanza de las Matemáticas
% Marzo-julio 2026, Grupo A
% ==============================================================================
% Uso recomendado:
% 1. Duplica este archivo por actividad y conserva el original como plantilla.
%    Ejemplo: PFM01_B1S1_Actividad3_SociedadesSigloXXI.tex
% 2. Actualiza el título visible de portada, la sección de actividad editable y la
%    fecha de entrega cuando corresponda.
% 3. Conserva las diapositivas de identidad PFM01 cuando la entrega requiera
%    contextualizar la unidad; si no son necesarias, comenta la sección completa.
% 4. No pegues la consigna completa en el producto final. Usa la consigna como
%    lista de cumplimiento y redacta una respuesta académica propia.
% 5. Toda actividad debe cerrar con postura profesional: posición, razón y efecto.
% 6. Si la actividad pide silueta, matriz, tabla, esquema, mapa, línea de tiempo,
%    infografía o producto visual, constrúyelo en LaTeX/TikZ cuando sea viable.
%
% Enfoque académico para PFM01:
% - Partir de la práctica docente como objeto de reflexión y mejora.
% - Relacionar enseñanza, aprendizaje, paradigmas, matemáticas y contexto social.
% - Evitar resúmenes planos: concepto -> función -> límite -> ejemplo de aula ->
%   implicación para la enseñanza de las matemáticas.
% - Integrar experiencia docente sin convertirla en anécdota suelta; usarla como
%   evidencia situada para analizar decisiones, creencias y necesidades.
% - Priorizar una voz profesional: clara, reflexiva, argumentada y situada.
%
% Estructura sugerida por actividad:
% [ ] Portada con unidad, bloque/sesión, actividad, alumno, matrícula y fecha.
% [ ] Objetivo de la actividad redactado con palabras propias.
% [ ] Problema o pregunta detonadora.
% [ ] Desarrollo con conceptos, criterios, experiencia, evidencias y análisis.
% [ ] Producto solicitado visible: silueta, matriz, tabla, esquema, etc.
% [ ] Aplicación a la práctica docente y a la enseñanza de las matemáticas.
% [ ] Conclusión con postura: posición, razón y consecuencia.
% [ ] Fuentes consultadas en formato consistente.
%
% Checklist antes de entregar:
% [ ] Título, actividad, bloque, sesión y fecha están actualizados.
% [ ] El objetivo coincide con la consigna del EVA.
% [ ] El producto solicitado aparece visible, rotulado y completo.
% [ ] La introducción plantea un problema o postura, no solo anuncia el tema.
% [ ] El desarrollo explica, conecta y analiza; no copia ni parafrasea sin aporte.
% [ ] Hay relación explícita con práctica docente y matemáticas.
% [ ] La conclusión incluye postura personal argumentada.
% [ ] Ortografía, acentos, nombres propios y fechas revisadas.
% [ ] El PDF compila sin errores y las imágenes institucionales cargan.
% ==============================================================================

\pdfminorversion=7
\documentclass[
	spanish,
	aspectratio=1610,
	xcolor={dvipsnames,table}
]{beamer}

% Lienzo 16:10 ampliado: da más espacio por diapositiva sin cambiar el diseño.
\geometry{paperwidth=19.2cm,paperheight=12cm}

% Idioma, tipografía y recursos visuales
\usepackage[utf8]{inputenc}
\usepackage[T1]{fontenc}
\usepackage[spanish,es-nodecimaldot]{babel}
\usepackage{lmodern}
\usepackage{booktabs}
\usepackage{graphicx}
\usepackage{ragged2e}
\usepackage{tikz}
\usepackage{hyperref}

% Datos del alumno y del programa
\newcommand{\studentname}{Mtro. Martin Jonathan de la Cruz Muñoz}
\newcommand{\studentshort}{M. de la Cruz Muñoz}
\newcommand{\studentid}{00874}
\newcommand{\studentemail}{martin.delacruz@campus.iiiepe.edu.mx}
\newcommand{\programname}{Maestría en Aprendizaje y Enseñanza de las Matemáticas}
\newcommand{\programshort}{MAEM}
\newcommand{\modality}{Modalidad presencial}
\newcommand{\periodname}{1er tetramestre marzo--julio 2026}
\newcommand{\coursecode}{PFM01}
\newcommand{\coursename}{Fundamentos para la enseñanza y el aprendizaje}
\newcommand{\coursegroup}{Grupo A}
\newcommand{\courseperiod}{Marzo--julio 2026}
\newcommand{\learningunit}{Unidad de aprendizaje}
\newcommand{\evaname}{EVA posgrado}
\newcommand{\blockname}{Bloque 1}
\newcommand{\sessionname}{Sesión 1}
\newcommand{\videoreflection}{https://youtu.be/WbOtm0zkxLQ?si=2mrILdsI5lxxwPju}
\newcommand{\scottreading}{El futuro del aprendizaje (I): ¿Por qué deben cambiar el contenido y los métodos de aprendizaje en el siglo XXI?}

% Datos institucionales
\newcommand{\institutionname}{IIIEPE}
\newcommand{\institutionlongname}{Instituto de Investigación, Innovación y Estudios de Posgrado para la Educación del Estado de Nuevo León}
\newcommand{\institutionurl}{https://iiiepe.edu.mx/}
\newcommand{\institutionplace}{Monterrey, Nuevo León}

% Recursos oficiales descargados desde el sitio del IIIEPE
\newcommand{\iiiepelogo}{img/iiiepe/logo-iiiepe-blanco.png}
\newcommand{\iiiepemark}{img/iiiepe/marca-iiiepe.png}
\newcommand{\iiiepecover}{img/iiiepe/portada-institucional.png}
\newcommand{\iiiepedots}{img/iiiepe/logo-puntos.png}
\newcommand{\maemplanimage}{img/iiiepe/maem-plan-estudios.png}
\newcommand{\maemsource}{https://iiiepe.edu.mx/maestria-en-aprendizaje-y-ensenanza-de-las-matematicas/}

% Paleta inspirada en la identidad visual del sitio institucional
\definecolor{iiiepeNavy}{HTML}{163B5C}
\definecolor{iiiepeBlue}{HTML}{1F6F9F}
\definecolor{iiiepeAqua}{HTML}{20B7A8}
\definecolor{iiiepeOrange}{HTML}{F28C28}
\definecolor{iiiepePaper}{HTML}{F6F8F7}
\definecolor{iiiepeInk}{HTML}{1E2A32}
\definecolor{iiiepeMuted}{HTML}{667085}

\mode<presentation>{
	\usetheme{default}
	\usefonttheme{professionalfonts}
	\setbeamertemplate{navigation symbols}{}
	\setbeamertemplate{blocks}[rounded][shadow=false]
	\setbeamercovered{invisible}
}

\hypersetup{
	colorlinks=true,
	urlcolor=iiiepeBlue,
	linkcolor=iiiepeNavy
}

\setbeamercolor{background canvas}{bg=white}
\setbeamercolor{normal text}{fg=iiiepeInk,bg=white}
\setbeamercolor{structure}{fg=iiiepeBlue}
\setbeamercolor{frametitle}{fg=white,bg=iiiepeNavy}
\setbeamercolor{title}{fg=white,bg=iiiepeNavy}
\setbeamercolor{subtitle}{fg=white}
\setbeamercolor{block title}{fg=white,bg=iiiepeBlue}
\setbeamercolor{block body}{fg=iiiepeInk,bg=iiiepePaper}
\setbeamercolor{alerted text}{fg=iiiepeOrange}
\setbeamerfont{title}{size=\LARGE,series=\bfseries}
\setbeamerfont{frametitle}{size=\large,series=\bfseries}
\setbeamerfont{block title}{series=\bfseries}
\setbeamertemplate{itemize item}{\textcolor{iiiepeAqua}{\large$\blacktriangleright$}}
\setbeamertemplate{itemize subitem}{\textcolor{iiiepeOrange}{\scriptsize$\blacksquare$}}
\setbeamertemplate{enumerate item}{\textcolor{iiiepeBlue}{\insertenumlabel.}}

\setbeamertemplate{frametitle}{
	\nointerlineskip
	\begin{beamercolorbox}[wd=\textwidth,ht=1.05cm,dp=0.22cm,leftskip=0.35cm,rightskip=0.25cm]{frametitle}
		\begin{minipage}[c]{0.72\textwidth}
			\insertframetitle
		\end{minipage}
		\hfill
		\raisebox{-0.1cm}{\includegraphics[height=0.58cm]{\iiiepelogo}}
	\end{beamercolorbox}
	\vspace{-0.04cm}
	{\color{iiiepeOrange}\rule{\textwidth}{1.4pt}}
}

\setbeamertemplate{footline}{
	\leavevmode%
	\hbox{%
		\begin{beamercolorbox}[wd=.34\paperwidth,ht=2.7ex,dp=1ex,leftskip=1em]{author in head/foot}%
			\color{white}\studentshort
		\end{beamercolorbox}%
		\begin{beamercolorbox}[wd=.38\paperwidth,ht=2.7ex,dp=1ex,center]{title in head/foot}%
			\color{white}\programshort\ · \modality
		\end{beamercolorbox}%
		\begin{beamercolorbox}[wd=.20\paperwidth,ht=2.7ex,dp=1ex,rightskip=1em plus 1fil]{date in head/foot}%
			\color{white}Matrícula \studentid\hfill\insertframenumber/\inserttotalframenumber
		\end{beamercolorbox}%
	}%
	\vskip0pt%
}
\setbeamercolor{author in head/foot}{bg=iiiepeNavy}
\setbeamercolor{title in head/foot}{bg=iiiepeBlue}
\setbeamercolor{date in head/foot}{bg=iiiepeNavy}

\newcommand{\accentbar}{%
	\begin{tikzpicture}[remember picture,overlay]
		\path[use as bounding box] (0,0) rectangle (0,0);
		\fill[iiiepeAqua] (current page.south west) rectangle ([xshift=0.42\paperwidth,yshift=0.08cm]current page.south west);
		\fill[iiiepeOrange] ([xshift=0.42\paperwidth]current page.south west) rectangle ([xshift=\paperwidth,yshift=0.08cm]current page.south west);
	\end{tikzpicture}
}

\newcommand{\sectionframe}[1]{%
	\begin{frame}[plain]
		\begin{tikzpicture}[remember picture,overlay]
			\fill[iiiepeNavy] (current page.south west) rectangle (current page.north east);
			\node[opacity=0.13,anchor=east] at ([xshift=0.8cm]current page.east) {\includegraphics[height=9.4cm]{\iiiepedots}};
			\fill[iiiepeAqua] ([xshift=0.85cm,yshift=2.1cm]current page.south west) rectangle ([xshift=1.05cm,yshift=6.3cm]current page.south west);
			\fill[iiiepeOrange] ([xshift=1.15cm,yshift=2.1cm]current page.south west) rectangle ([xshift=1.35cm,yshift=5.35cm]current page.south west);
		\end{tikzpicture}
		\vspace{1.8cm}
		\hspace{1.25cm}
		\begin{minipage}{0.72\paperwidth}
			{\color{white}\Huge\bfseries #1}\\[0.25cm]
			{\color{white!75}\programname}
		\end{minipage}
	\end{frame}
}

\AtBeginSection[]{
	\sectionframe{\insertsectionhead}
}

\title[\coursecode]{\coursename}
\subtitle{\programname}
\author[\studentshort]{\texorpdfstring{\studentname\\\footnotesize Matrícula: \studentid\\\href{mailto:\studentemail}{\studentemail}}{\studentname, Matricula: \studentid}}
\institute[\institutionname]{\institutionlongname\\\coursecode\ · \coursegroup}
\date[\courseperiod]{\institutionplace\\\courseperiod}

\begin{document}

% Portada institucional
\begin{frame}[plain]
	\begin{tikzpicture}[remember picture,overlay]
		\node[anchor=center,opacity=0.18] at (current page.center) {\includegraphics[width=\paperwidth,height=\paperheight]{\iiiepecover}};
		\fill[iiiepeNavy,opacity=0.94] (current page.south west) rectangle ([xshift=0.55\paperwidth]current page.north west);
		\fill[iiiepeBlue,opacity=0.92] ([xshift=0.55\paperwidth]current page.south west) rectangle (current page.north east);
		\node[anchor=north west] at ([xshift=0.55cm,yshift=-0.45cm]current page.north west) {\includegraphics[width=4.15cm]{\iiiepelogo}};
		\node[opacity=0.17,anchor=south east] at ([xshift=0.9cm,yshift=-0.3cm]current page.south east) {\includegraphics[height=8.8cm]{\iiiepedots}};
		\fill[iiiepeOrange] ([xshift=0.56\paperwidth,yshift=1.15cm]current page.south west) rectangle ([xshift=0.92\paperwidth,yshift=1.28cm]current page.south west);
		\fill[iiiepeAqua] ([xshift=0.56\paperwidth,yshift=0.88cm]current page.south west) rectangle ([xshift=0.82\paperwidth,yshift=1.01cm]current page.south west);
	\end{tikzpicture}

	\vspace{1.9cm}
	\begin{minipage}{0.50\paperwidth}
		{\color{white}\Huge\bfseries Fundamentos para la enseñanza\\ y el aprendizaje}\\[0.25cm]
		{\color{white!86}\large \coursecode\ · \coursegroup}\\[0.35cm]
		{\color{iiiepeOrange}\rule{0.82\linewidth}{1.3pt}}\\[0.45cm]
		{\color{white}\studentname}\\
		{\color{white!80}\footnotesize Matrícula: \studentid\ · \studentemail}
	\end{minipage}
	\hfill
	\begin{minipage}{0.36\paperwidth}
		\raggedleft
		{\color{white}\Large\bfseries \institutionname}\\[0.18cm]
		{\color{white!82}\footnotesize \institutionlongname}\\[0.5cm]
		{\color{white}\small \programname}\\
		{\color{white!78}\small \courseperiod}\\[0.5cm]
		{\color{white!74}\footnotesize \institutionplace\ · \learningunit}\\
		{\color{white!74}\footnotesize \href{\institutionurl}{\institutionurl}}
	\end{minipage}
\end{frame}

\begin{frame}{Contenido}
	\tableofcontents
\end{frame}

\section{Identidad PFM01}

\begin{frame}{Ficha de la unidad}
	\begin{columns}[T]
		\begin{column}{0.42\textwidth}
			\begin{center}
				\includegraphics[width=0.54\linewidth]{\iiiepemark}\\[0.35cm]
				{\Large\bfseries\color{iiiepeNavy}\institutionname}\\
				{\small\color{iiiepeMuted}\evaname}
			\end{center}
		\end{column}
		\begin{column}{0.54\textwidth}
			\begin{block}{Alumno}
				\textbf{\studentname}\\
				Matrícula: \studentid\\
				Correo: \href{mailto:\studentemail}{\studentemail}
			\end{block}
			\begin{block}{Unidad de aprendizaje}
				\textbf{\coursecode\ · \coursename}\\
				\coursegroup\\
				\courseperiod
			\end{block}
		\end{column}
	\end{columns}
	\accentbar
\end{frame}

\begin{frame}{Contexto académico}
	\begin{columns}[T]
		\begin{column}{0.48\textwidth}
			\begin{block}{Datos generales}
				\begin{itemize}
					\item Programa: \programname.
					\item Unidad: \coursename.
					\item Clave: \coursecode.
					\item Periodo: \courseperiod.
					\item Grupo: \coursegroup.
				\end{itemize}
			\end{block}
		\end{column}
		\begin{column}{0.48\textwidth}
			\begin{block}{Propósito de la unidad}
				Analizar los fundamentos de la enseñanza y el aprendizaje en el
				contexto actual, reconociendo el papel docente, las características
				de las clases de matemáticas y las habilidades necesarias para
				promover aprendizajes significativos.
			\end{block}
		\end{column}
	\end{columns}
	\accentbar
\end{frame}

\begin{frame}{Introducción a la unidad}
	\begin{block}{Sentido de la unidad}
		La unidad analiza el papel del docente en las sociedades del siglo XXI,
		los cambios educativos actuales y las implicaciones que estos tienen en
		la enseñanza de las matemáticas.
	\end{block}
	\begin{columns}[T]
		\begin{column}{0.32\textwidth}
			\begin{block}{Docencia}
				Reconocer identidad, decisiones, creencias y áreas de fortalecimiento.
			\end{block}
		\end{column}
		\begin{column}{0.32\textwidth}
			\begin{block}{Aprendizaje}
				Revisar tipos de aprendizaje, habilidades y condiciones actuales.
			\end{block}
		\end{column}
		\begin{column}{0.32\textwidth}
			\begin{block}{Matemáticas}
				Vincular los fundamentos educativos con la práctica matemática escolar.
			\end{block}
		\end{column}
	\end{columns}
	\accentbar
\end{frame}

\begin{frame}{Temario base de la unidad}
	\begin{enumerate}
		\item Introducción: enseñanza, aprendizaje, paradigmas y matemáticas.
		\item Sociedades del siglo XXI y transformación educativa.
		\item Tipos de aprendizaje en las sociedades actuales.
		\item Siete habilidades básicas para el contexto contemporáneo.
		\item Sugerencias metodológicas para fortalecer la práctica educativa.
		\item Enseñanza de las matemáticas y características de sus clases.
		\item Programa 2022 de la Nueva Escuela Mexicana.
		\item Trabajo final integrador.
	\end{enumerate}
	\accentbar
\end{frame}

\section{\blockname\ · \sessionname}

\begin{frame}{Introducción a la sesión}
	\begin{block}{Punto de partida}
		La primera sesión promueve una reflexión inicial sobre conceptos
		fundamentales del campo educativo: enseñanza, aprendizaje, paradigmas y
		matemáticas.
	\end{block}
	\begin{itemize}
		\item Recuperar experiencia profesional y saberes previos.
		\item Identificar creencias y perspectivas sobre la práctica docente.
		\item Construir significados desde el diálogo y el análisis crítico.
		\item Evitar respuestas únicas; privilegiar reflexión argumentada.
	\end{itemize}
	\accentbar
\end{frame}

\begin{frame}{Recurso de reflexión}
	\begin{columns}[T]
		\begin{column}{0.48\textwidth}
			\begin{block}{Video}
				El recurso invita a reflexionar sobre paradigmas educativos
				tradicionales y la necesidad de replantear la enseñanza y el
				aprendizaje en las sociedades actuales.
			\end{block}
		\end{column}
		\begin{column}{0.48\textwidth}
			\begin{block}{Clave de análisis}
				Relacionar el planteamiento del video con la experiencia docente y
				con la enseñanza de las matemáticas.
			\end{block}
			{\scriptsize \href{\videoreflection}{\videoreflection}}
		\end{column}
	\end{columns}
	\accentbar
\end{frame}

\begin{frame}{Actividad diagnóstica}
	\begin{block}{Propósito}
		Identificar ideas previas sobre enseñanza, aprendizaje, paradigmas y
		matemáticas para establecer un punto de partida de la unidad.
	\end{block}
	\scriptsize
	\begin{tabular}{p{0.20\textwidth}p{0.70\textwidth}}
		\toprule
		\textbf{Concepto} & \textbf{Pregunta orientadora} \\
		\midrule
		Enseñanza & ¿Qué decisiones tomo para orientar, mediar o transformar el aprendizaje? \\
		Aprendizaje & ¿Cómo reconozco que un estudiante construyó comprensión matemática? \\
		Paradigmas & ¿Qué creencias sostienen mi forma de enseñar y evaluar? \\
		Matemáticas & ¿Las presento como procedimiento, lenguaje, razonamiento o herramienta cultural? \\
		\bottomrule
	\end{tabular}
	\accentbar
\end{frame}

\begin{frame}{Actividad: ¿Cómo soy como maestro?}
	\begin{block}{Propósito}
		Analizar la identidad docente, reconocer fortalezas, áreas de oportunidad
		y necesidades de formación para fortalecer la práctica profesional en la
		enseñanza de las matemáticas.
	\end{block}
	\begin{itemize}
		\item Elaborar una silueta que represente la identidad docente.
		\item Explicar fortalezas y áreas de oportunidad desde la práctica real.
		\item Identificar bibliografía o formación necesaria para mejorar.
		\item Compartir hallazgos en diálogo presencial y reflexión grupal.
	\end{itemize}
	\accentbar
\end{frame}

\begin{frame}{Silueta docente: organización}
	\begin{columns}[T]
		\begin{column}{0.32\textwidth}
			\begin{block}{Lado izquierdo}
				Fortalezas y áreas de oportunidad como docente de matemáticas,
				explicadas con ejemplos de práctica.
			\end{block}
		\end{column}
		\begin{column}{0.32\textwidth}
			\begin{block}{Centro}
				Definición personal como docente y rasgos que caracterizan la
				práctica educativa.
			\end{block}
		\end{column}
		\begin{column}{0.32\textwidth}
			\begin{block}{Lado derecho}
				Necesidades de formación o bibliografía justificadas por su aporte
				al desarrollo profesional.
			\end{block}
		\end{column}
	\end{columns}
	\accentbar
\end{frame}

\section{Actividad 3}

\begin{frame}{Sociedades del siglo XXI}
	\begin{block}{Propósito de la actividad}
		Analizar los factores que impulsan la transformación de la enseñanza y
		construir una postura fundamentada sobre su impacto en la educación y en
		la enseñanza de las matemáticas.
	\end{block}
	\begin{itemize}
		\item Lectura base: \textit{\scottreading}, de Cynthia Luna Scott.
		\item Apertura: sábado 25 de abril de 2026, 8:01 a. m.
		\item Entrega: viernes 1 de mayo de 2026, 11:59 p. m.
	\end{itemize}
	\accentbar
\end{frame}

\begin{frame}{Factores de transformación}
	\scriptsize
	\begin{columns}[T]
		\begin{column}{0.48\textwidth}
			\begin{itemize}
				\item Competencias nuevas para un mundo complejo.
				\item Nuevas características de las y los estudiantes.
				\item Desinterés juvenil y abandono prematuro.
				\item Valor actual de la educación formal.
				\item Escasez de preparación.
				\item Evolución del mercado de trabajo.
			\end{itemize}
		\end{column}
		\begin{column}{0.48\textwidth}
			\begin{itemize}
				\item Cambios en modelos de enseñanza y aprendizaje.
				\item Nuevos medios didácticos y vías de aprendizaje.
				\item Nuevas normas de evaluación y rendición de cuentas.
				\item Aprendizaje en cualquier momento y lugar.
				\item Crecimiento acelerado de la información.
				\item Del alumno consumidor al productor de conocimiento.
			\end{itemize}
		\end{column}
	\end{columns}
	\accentbar
\end{frame}

\begin{frame}{Matriz para argumentar cada factor}
	\scriptsize
	\begin{tabular}{p{0.22\textwidth}p{0.22\textwidth}p{0.22\textwidth}p{0.20\textwidth}}
		\toprule
		\textbf{Factor} & \textbf{Relevancia} & \textbf{Relación educativa} & \textbf{Postura docente} \\
		\midrule
		Escribir factor & Explicar por qué transforma la enseñanza & Vincular con aula, estudiantes o currículo & Sostener posición con razón y efecto \\
		\bottomrule
	\end{tabular}
	\vspace{0.45cm}
	\begin{block}{Regla de escritura}
		No limitarse a describir: cada factor debe convertirse en argumento,
		con relevancia, relación con el contexto actual y conexión con la
		enseñanza de las matemáticas.
	\end{block}
	\accentbar
\end{frame}

\begin{frame}{Recursos de apoyo}
	\begin{block}{Lectura principal}
		Cynthia Luna Scott, \textit{\scottreading}. Usar la lectura para
		construir argumentos sobre cambio educativo, no para copiar definiciones.
	\end{block}
	\begin{block}{Organizador visual}
		La infografía de procesos y etapas puede servir como referencia de forma:
		jerarquizar, ordenar, sintetizar y conectar ideas dentro de productos
		visuales propios.
	\end{block}
	\begin{block}{Criterio de uso}
		Todo recurso debe convertirse en análisis situado: qué plantea, por qué
		importa, cómo se relaciona con el aula y qué exige a la enseñanza de las
		matemáticas.
	\end{block}
	\accentbar
\end{frame}

\section{Pautas de plantilla}

\begin{frame}{Ruta sugerida para cada actividad}
	\begin{enumerate}
		\item Leer la consigna y convertirla en objetivo propio.
		\item Identificar producto esperado: silueta, matriz, exposición, análisis o propuesta.
		\item Separar conceptos, criterios, evidencias y experiencia docente.
		\item Construir argumento: problema, fundamento, aplicación y consecuencia.
		\item Integrar relación explícita con la enseñanza de las matemáticas.
		\item Cerrar con postura profesional: posición, razón y efecto.
	\end{enumerate}
	\accentbar
\end{frame}

\begin{frame}{Criterios de calidad para entregar}
	\begin{itemize}
		\item Claridad: cada diapositiva debe sostener una idea principal.
		\item Pertinencia: responder a la actividad, no solo al tema general.
		\item Evidencia: incluir fuentes, ejemplos, casos o productos verificables.
		\item Análisis: explicar relaciones, límites, implicaciones y decisiones didácticas.
		\item Identidad profesional: conectar aprendizaje, experiencia y mejora docente.
		\item Escritura propia: evitar copiar o parafrasear sin argumento.
	\end{itemize}
	\accentbar
\end{frame}

\begin{frame}{Checklist antes de entregar}
	\begin{itemize}
		\item Título, unidad, bloque, sesión, alumno y fecha actualizados.
		\item Objetivo alineado con la consigna real del EVA.
		\item Producto solicitado visible, rotulado y completo.
		\item Relación con práctica docente y enseñanza de matemáticas.
		\item Conclusión con postura personal argumentada.
		\item Fuentes revisadas y citadas de forma consistente.
	\end{itemize}
	\accentbar
\end{frame}

\section{Plantilla editable}

\begin{frame}{Título de la actividad}
	\begin{columns}[T]
		\begin{column}{0.62\textwidth}
			\begin{itemize}
				\item Sustituir este punto por la idea principal de la actividad.
				\item Agregar evidencias, argumentos o referencias de clase.
				\item Cerrar con una implicación para la práctica docente.
			\end{itemize}
		\end{column}
		\begin{column}{0.32\textwidth}
			\begin{block}{Enfoque}
				Fundamentos de enseñanza y aprendizaje aplicados a la práctica
				docente en matemáticas.
			\end{block}
		\end{column}
	\end{columns}
	\accentbar
\end{frame}

\begin{frame}{Objetivo de la actividad}
	\begin{block}{Redactar objetivo}
		Escribir aquí el objetivo de la actividad en una frase propia, usando un
		verbo observable: analizar, comparar, diseñar, evaluar, explicar,
		proponer, argumentar o sintetizar.
	\end{block}
	\begin{itemize}
		\item ¿Qué se va a estudiar o producir?
		\item ¿Con qué criterio teórico o didáctico se trabajará?
		\item ¿Para qué mejora, decisión o aprendizaje servirá?
	\end{itemize}
	\accentbar
\end{frame}

\begin{frame}{Análisis}
	\begin{columns}[T]
		\begin{column}{0.48\textwidth}
			\begin{block}{Concepto clave}
				Escribir aquí la definición, postura teórica o problema central.
			\end{block}
		\end{column}
		\begin{column}{0.48\textwidth}
			\begin{block}{Aplicación}
				Relacionar el concepto con la enseñanza y aprendizaje de las matemáticas.
			\end{block}
		\end{column}
	\end{columns}
	\vfill
	\begin{center}
		{\color{iiiepeMuted}\small \institutionname\ · \coursecode\ · \courseperiod}
	\end{center}
	\accentbar
\end{frame}

\begin{frame}{Producto solicitado}
	\begin{block}{Tipo de producto}
		Indicar si la actividad pide cuadro comparativo, mapa conceptual, línea
		de tiempo, exposición, ensayo breve, análisis de caso, propuesta
		didáctica, secuencia, rúbrica o ambiente de aprendizaje.
	\end{block}
	\begin{block}{Evidencia de elaboración}
		Insertar aquí el producto construido o una síntesis visual editable.
	\end{block}
	\accentbar
\end{frame}

\begin{frame}{Aplicación a la práctica docente}
	\begin{itemize}
		\item Contexto educativo o grupo al que podría aplicarse.
		\item Problema matemático, didáctico o evaluativo que atiende.
		\item Estrategia de intervención o mejora propuesta.
		\item Evidencia esperada para valorar su impacto.
	\end{itemize}
	\accentbar
\end{frame}

\section{Cierre}

\begin{frame}{Conclusiones}
	\begin{enumerate}
		\item Sintetizar el aprendizaje principal.
		\item Explicar su relevancia para el contexto educativo.
		\item Plantear una mejora o pregunta para continuar investigando.
	\end{enumerate}
	\accentbar
\end{frame}

\begin{frame}[plain]
	\begin{tikzpicture}[remember picture,overlay]
		\fill[iiiepeNavy] (current page.south west) rectangle (current page.north east);
		\node[opacity=0.13,anchor=south east] at ([xshift=0.8cm,yshift=-0.2cm]current page.south east) {\includegraphics[height=9cm]{\iiiepedots}};
		\node[anchor=north west] at ([xshift=0.55cm,yshift=-0.45cm]current page.north west) {\includegraphics[width=3.7cm]{\iiiepelogo}};
	\end{tikzpicture}
	\vspace{2.2cm}
	\begin{center}
		{\color{white}\Huge\bfseries Gracias}\\[0.55cm]
		{\color{white}\Large \studentname}\\[0.18cm]
		{\color{white!78}\footnotesize Matrícula \studentid\ · \studentemail}\\[0.45cm]
		{\color{iiiepeOrange}\rule{0.36\paperwidth}{1.4pt}}\\[0.35cm]
		{\color{white!72}\footnotesize \institutionname\ · \coursecode\ · \coursename}
	\end{center}
\end{frame}

\end{document}
